\documentclass{article}

% Pakete
\usepackage[utf8]{inputenc}
\usepackage[T1]{fontenc}
\usepackage[german]{babel}
\usepackage{amsmath, amssymb}
\usepackage{graphicx}
\usepackage{hyperref}
\usepackage{geometry}
\geometry{margin=2.5cm}

\title{Einführung in LaTeX}
\author{Dein Name}
\date{\today}

\begin{document}

\maketitle
\tableofcontents
\newpage

\section{Was ist LaTeX?}
LaTeX ist ein Textsatzsystem, das besonders für wissenschaftliche Dokumente verwendet wird.
Es eignet sich hervorragend für:
\begin{itemize}
    \item mathematische Formeln
    \item strukturierte Dokumente
    \item wissenschaftliche Arbeiten
\end{itemize}

Im Gegensatz zu Word arbeitet LaTeX nicht visuell, sondern mit Code.

\section{Grundstruktur eines Dokuments}
Jedes LaTeX-Dokument besteht aus zwei Teilen:
\begin{enumerate}
    \item Präambel (alles vor \texttt{\textbackslash begin\{document\}})
    \item Dokumentinhalt
\end{enumerate}

\subsection{Beispiel}
\begin{verbatim}
\documentclass{article}
\begin{document}
Hallo Welt!
\end{document}
\end{verbatim}

\section{Textformatierung}
\subsection{Fett, kursiv, unterstrichen}
\textbf{Fetter Text}

\textit{Kursiver Text}

\underline{Unterstrichener Text}

\subsection{Absätze}
Ein neuer Absatz entsteht durch eine Leerzeile.

Dies ist ein neuer Absatz.

\section{Mathematik in LaTeX}
LaTeX ist besonders stark bei mathematischen Formeln.

\subsection{Inline-Mathe}
Dies ist eine Formel im Text: $a^2 + b^2 = c^2$

\subsection{Abgesetzte Formeln}
\[
E = mc^2
\]

\subsection{Brüche und Wurzeln}
\[
\frac{a}{b}, \quad \sqrt{x}, \quad \sqrt[n]{x}
\]

\subsection{Summen und Integrale}
\[
\sum_{i=1}^{n^2} i^2
\]

\[
\int_{0}^{1} x^2 \, dx
\]

\section{Listen}
\subsection{Aufzählung}
\begin{itemize}
    \item Punkt 1
    \item Punkt 2
    \item Punkt 3
\end{itemize}

\subsection{Nummerierte Liste}
\begin{enumerate}
    \item Erster Punkt
    \item Zweiter Punkt
\end{enumerate}

\section{Tabellen}
\begin{tabular}{|c|c|c|}
\hline
Spalte 1 & Spalte 2 & Spalte 3 \\
\hline
1 & 2 & 3 \\
4 & 5 & 6 \\
\hline
\end{tabular}

\section{Bilder einfügen}
\begin{figure}[h]
\centering
\caption{Ein Beispielbild}
\end{figure}

\section{Referenzen und Links}
Ein Beispiel-Link:
\url{https://www.latex-project.org}

\section{Fazit}
LaTeX ist ein mächtiges Werkzeug für professionelle Dokumente. Anfangs wirkt es komplex, aber mit etwas Übung ist es sehr effizient.

\end{document}